\section{Verification purpose}

This case verifies raster lookup and fire-behavior coupling when a propagating front crosses four fuel-model regions.  Flat terrain and zero wind isolate fuel identity.  The verification question is whether every quadrant's maximum spread rate is within 10 percent of its fuel-specific reference.

The scientific setup and historical expected response are reconstructed from Section~3.2 of the supplied \emph{ELMFIRE Guide}. The Guide is used as reference provenance, not as an instruction or a substitute for evidence. Configuration statements below are checked against the generated GeoTIFFs and namelist; current results and decisions come only from case-specific outputs and generated metrics. Where the generated experiment differs from the historical example, the report states that difference explicitly.

\section{Mathematical formulation and numerical implementation}

The level-set normal speed is evaluated from the fuel model at each active cell.  Because the environmental forcing is uniform, changes in \(V_n\) across the domain should follow only changes in the fuel-model classification raster.  Correct row orientation and quadrant masking are therefore part of the test.

At each active cell the fuel code selects a complete Rothermel parameter row, and the local no-wind/no-slope rate is
\begin{equation}
R_0(x,y)=\frac{I_R(x,y)\,\xi(x,y)}{\rho_b(x,y)\,\epsilon(x,y)\,Q_{ig}(x,y)}.
\end{equation}
Thus the raster lookup changes the full nonlinear fuel response, not a post-hoc scalar multiplier. The formulation follows Andrews (2018); the Guide supplies the four BEHAVE reference values.

\section{Derivation of expected behavior}

The independent regional targets are \(0.080\), \(0.472\), \(1.492\), and \(0.871~\mathrm{m\,min^{-1}}\) for Fuel Models 8, 7, 4, and 2 in the northwest, northeast, southwest, and southeast quadrants, respectively.  These values also predict a visibly noncircular front after it enters all four regions. The Guide historically reported ELMFIRE rates of 0.080, 0.551, 1.974, and 0.872 m/min for FB8, FB7, FB4, and FB2, noting discrepancies for FB7 and FB4. Those historical outputs are preserved here as context only; the generated metric table below is the current evidence and is never replaced by those values.

% BEGIN EXPLICIT METRIC DERIVATION
All four references use the no-wind/no-slope Rothermel equation
\(R_0=I_R\xi/(\rho_b\epsilon Q_{ig})\) with
\(Q_{ig}=250+1116M_f\) and the generated moisture vector
\((0.06,0.07,0.08,0.60,0.90)\).  Substitution of each generated standard
fuel model gives
\[
\begin{array}{c|cc}
\text{fuel model}&I_R~(\mathrm{kW\,m^{-2}})&R_0~(\mathrm{m\,min^{-1}})\\
\hline
8&175.566&0.079552\\
7&406.479&0.470579\\
4&2272.300&1.493256\\
2&669.483&0.871756
\end{array}
\]
and hence the report references \(0.080,\ 0.472,\ 1.492,\) and
\(0.871~\mathrm{m\,min^{-1}}\), respectively.  For each fuel model,
\(e_R=|R_{\max}-R_{\mathrm{ref}}|/R_{\mathrm{ref}}\); the table substitutes
the corresponding row above and the prescribed relative-error limit.
% END EXPLICIT METRIC DERIVATION

\section{Verification metrics and acceptance criteria}

The maximum finite spread rate in each quadrant is compared with its corresponding target.  Regional maxima preserve the local head-fire response and avoid mixing physically distinct fuels.  Every component uses 10 percent relative tolerance, and PASS requires all four components.

\section{Simulation configuration and predicted outcome}

Four equal quadrants on a \(126\times126\), 30 m grid contain Fuel Models 8, 7, 4, and 2. Terrain is flat, wind is zero, moisture is uniform, and ignition is central. The fastest extension is predicted in Fuel Model 4 and the slowest in Fuel Model 8.

\begin{center}
\begin{tabular}{@{}p{0.20\textwidth}p{0.20\textwidth}p{0.12\textwidth}p{0.38\textwidth}@{}}
\toprule
Parameter & Configured value & Units & Role in the experiment \\
\midrule
Grid & \(126\times126\) at 30 & m per cell & Resolves the whole spatial response on the generated north-up raster \\
Projection & EPSG:32635 & --- & Fixes the map-coordinate and row-orientation convention \\
Time control & \(\Delta t=1\), \(t_{stop}=259200\) & s & Whole-second step; stop time is an exact timestep multiple \\
Meteorology interval & 3600 & s & Defines weather-band timing and stalled-final compatibility \\
Live moisture & herb 60; woody 90 & percent & Holds live-fuel damping fixed while the named mechanism is isolated \\
Evidence & prepared inputs; final case outputs & --- & Separates configuration provenance from acceptance evidence \\
\bottomrule
\end{tabular}
\end{center}

Figure~\ref{fig:input-configuration} shows the prescribed spatial condition used in this experiment. The displayed field represents the initial level-set geometry or a controlling fuel, terrain, moisture, wind, or canopy quantity, as identified by the figure labels. The plotted values come from the prepared simulation inputs, not from the expected response.

\begin{figure}[H]
  \centering
  \IfFileExists{../figures/input_configuration.pdf}{%
    \includegraphics[width=0.98\textwidth,height=0.46\textheight,keepaspectratio]{input_configuration.pdf}%
  }{\fbox{\parbox{0.88\textwidth}{Prepared-input evidence is unavailable. The prescribed configuration is described in the text; this absence does not establish that it was simulated.}}}
  \caption{Prepared input conditions from the prescribed spatial fields. Coordinates are in metres; the colorbar identifies the displayed field and its physical units or category. The ignition panel is derived directly from \(\phi\leq0\).}
  \label{fig:input-configuration}
\end{figure}

\section{Actual simulation results}

The required evidence is generated only from case-specific ELMFIRE outputs. The current postprocessing status is \textbf{\VerificationStatus}; \PassingMetricCount{} of \MetricCount{} required metric rows pass across \VariantCount{} simulation condition(s). A missing output is reported as NOT RUN or NOT EVALUABLE and is not interpreted as verification evidence.

\begin{figure}[H]
  \centering
  \IfFileExists{../figures/whole_domain_result.pdf}{%
    \includegraphics[width=0.82\textwidth,height=0.58\textheight,keepaspectratio]{whole_domain_result.pdf}%
  }{\fbox{\parbox{0.85\textwidth}{No case-specific whole-domain ELMFIRE raster is available.}}}
  \caption{Final whole-domain ELMFIRE time of arrival for the displayed simulation condition, in seconds. The coordinate reference system is stated in the figure. The calculated metrics, not visual appearance, determine acceptance.}
\end{figure}

\section{Calculated metrics and verification decision}

\begin{center}
\begin{tabular}{@{}llll@{}}
\toprule
Simulation condition / metric & Expected & Calculated & Status \\
\midrule
\MetricRows
\bottomrule
\end{tabular}
\end{center}

The displayed overall decision is \textbf{\VerificationStatus}. When calculated verification and validation metrics exists, result analysis procedure generates every table value and status; when it is absent, the report states NOT EVALUABLE with no invented rows.

\section{Technical interpretation}

A quadrant swap indicates raster orientation or mask errors; a fuel-specific discrepancy indicates fuel-table selection or fire-behavior coupling. No scientific conclusion should be drawn when the status is NOT RUN or NOT EVALUABLE.

\section{Scope and limitations}

The test samples four fuel models under one moisture state and does not validate the empirical fuel parameters. This case provides implementation verification evidence only.
